\section{Introduction}
\begin{table}[h]
    \centering
    \begin{tabular}{l}
        \textit{Tra le rossastre nubi}                 \\
        \textit{Stormi d’uccelli neri,}                \\
        \textit{Com’esuli pensieri,}                   \\
        \textit{Nel vespero migrar.}                   \\
        \noalign{\smallskip}
        \scriptsize Among the reddish hazes,           \\
        \scriptsize flocking birds as black as night,  \\
        \scriptsize indistinct as thoughts in flight,  \\
        \scriptsize t’wards the setting sun they flee. \\
        \noalign{\smallskip}
        \scriptsize -- G. Carducci (Translation by M. Colarossi)
    \end{tabular}
\end{table}
The behaviour of flocking birds has been a subject of fascination for the human race for as long as is remembered. The study of such behaviour in a scientific way is much more novel; a flock of birds is studied as active matter \cite{vicsek_review}, and as such uses the techniques of the field, which is a novel one, since the study of such things has been limited by computational capability and other material constraints. Speaking of material constraints, for much of the history of the field there was no available data to study for the phenomenon of collective motion of birds. The first model to simulate it, called "boids", was created in 1987 by a computer graphic researcher specialized in movies and videogames, and was based on reasonable assumptions and qualitative evaluation \cite{reynolds} \cite{boids}.  \\
In the following years, the field evolved towards a more scientific approach. Jump started by T. Vicsek with his work on active matter (which was later disputed \cite{chate}), which was a simplified "boids" model created for the simulation of self propelled particles \cite{vicsek} \cite{vicsek_manual} \cite{vicsek_review}, the study of collective motion saw a rapid development with the creation of many models; first, Y. Tu and J. Toner applied fluid dynamics to the system, obtaining a rigorous theoretical model based on the Navier-Stokes equation \cite{toner_tu}. Then, I.D. Couzin developed the so-called "zonal" model, which was a refined Reynolds model \cite{couzin}. But alas, as scientific as those works might have been, they were still based on reasonable assumptions and qualitative evaluation, or, when it was quantitative, it was still testing reasonable quantities, not empirically verified ones. \\
That abruptly changed when, in 2008, A. Cavagna and his group published the first works about the STARFLAG system \cite{cavagna2008starflag1} \cite{cavagna2008starflag2} which set out, successfully, with building a comprehensive dataset on collective animal behaviour; in particular, they recorded a great number of flocks of the common starling (\textit{Sturnus vulgaris}) and their velocity, finally building the dataset the field needed. This group in particular was at the forefront of the field, first in analyzing their own dataset and finding the first emergent quantity of flocks, namely the shape of the so-called correlation function \cite{cavagna_2010}, then by analyzing it using statistical mechanics and by consequently using the MaxEnt statistical model to prove how the large scale interactions in a flock are determined by the small scale ones \cite{cavagna_2012}, and then by building a predictive dynamical model akin to a Hamiltonian formulation for rotations \cite{cavagna_2014}. In the meantime, their dataset sparked interest from other groups, thus leading to other dynamical models such as H. Hildenbrandt's one, which is a refined "zonal" model \cite{h_and_h}, in turn influenced by M. Bellerini's finding that birds do not follow metric rules but topological ones, namely they react only to the 7 or 6 nearest flockmates \cite{bellerini}. \\
In this work, we test the pre-2010 models in light of the findings by the Rome group, in particular that the shape of the correlation function is an emergent property of flocks \cite{cavagna_2010}. To be more precise, we implement Reynolds', Vicsek's and Couzin's models, with predator and obstacle avoidance, and we test the similarity between their correlation function and the empirical one. The rationale for choosing the aforementioned models is that the post-2010 ones have already been tested in their respective seminal papers \cite{cavagna_2014} \cite{h_and_h}, and the Toner-Tu model has technically already been tested by A. Cavagna, since his 2014 model reduces to Toner-Tu's formulation for large enough flocks \cite{cavagna_2014}. Finally, Chaté's model \cite{chate}, which is the one created when disputing Vicsek's one, was tested by Cavagna's group \cite{cavagna_2012}.